\documentclass[12pt]{article}
\usepackage[utf8]{inputenc}
\usepackage{setspace}
\usepackage[margin=1in]{geometry}
\usepackage{parskip}
\setlength{\parindent}{0pt}

\title{Sitting with Difficulty on the Instrumentalization of Religion for a Modernist Agenda}
\author{Hassan Akhter (ha10609)}
\date{}

\begin{document}
\maketitle

\section*{The Instrumentalization of Religion to achieve political ends and its philosophical reconciliation and reframing}
Throughout the course, I was at odds and in conflict with the current role of Islam in the modern world. Modern violence, the deepening of differences, and the shift away from the healing message of God toward modernist political and nationalist agendas have mutated religion into something which I have come to believe was never meant to be. Islam, and more so religion in general, has increasingly become a means to achieve political ends. Political entities deliberately use religion to wage division and conflict, using it as a tool to achieve their goals. Religion has been increasingly instrumentalized over the years to justify hate, division, and violence, with states using it as policy justification to wage war and secure political domination over others. I want to explore how this came to be and how we can reconcile this modernist reshaping of religion and find a way back to a more healing understanding of Islam and Religion in general. An understanding that does not immediately associate faith with divisive politics and violence.

This instrumentalization of religion is perhaps most evident in the formation of modern religious identities; people box themselves into their own bubbles of like-minded people and as a result presenting a homogeneous identity. In modern times this "collapse of the transcendent into the imminent" is more pronounced than ever (Naqvi, 2024). Religion is stripped of its transcendent, spiritual essence and reduced to a mere utility. This utility is attached to the divine to serve the interests of actors who seek to achieve their own ends with this association. As political entities construct modern national identities, they co-opt religious symbols, rhetoric, and morals to legitimize their power. As a result, politics actively distorts religion, shaping it to fit the narrow confines of nationalist, and modernist agendas, while religion, in turn, is used to distort political discourse, shutting down rational debate through claims of divine law.

\vspace{2em}

However, all these claims about political instrumentalization might lead us to believe that all instrumentalization is bad, whereas it could also be illustrated how religion is used to do good in society. All the ways in which it compels people to be accountable and act in a way that contributes to society, above our own self-interest, for reference the religiously ordained act of Zakat, it would be hard to argue how this use of religion to redistribute wealth in the society is not a net positive. In fact, when all religions first began they brought a message that was at the time very progressive and such a radical instrument of change to their way of life that it brought about persecution to the Prophets who preached it.

Conversely, to say that this weaponization of faith is a modern invention would be disingenuous, as it runs into a well-documented historical complication: hasn't this always been the case? We must acknowledge that religion and politics have almost always been intertwined. Throughout history, religion has frequently served as a tool for conquest and influence. Examples abound, from the sweeping expansions during the Rashidun and Umayyad Caliphates to the holy wars of the Crusades. In these eras, the banner of faith was linked to the expansion of empires and the consolidation of political power.

\vspace{2em}

To understand this, we must look at how empires actively used politics to fracture and divide religious communities for their own consolidation of power. In the early Islamic expansions, particularly during the transition into the Umayyad Caliphate, what was a unified faith community was deeply divided by political struggles over succession and leadership. The “Sunni/Shia” divide in Islam which persists today was not initially born of theological differences, but rather political conflicts over the right to rule. The Umayyad dynasty, for example, used religious rhetoric to legitimize its dynastic control and suppress political dissent, using state power to permanently fracture the unity of the early Muslims, a conflict epitomized by the tragedy of Karbala (Associated Press of Pakistan, n.d.). Similarly, during the Crusades, European political and religious leaders utilized the banner of holy war to politically unify a fractured, warring Europe by creating a vilified 'other' (Falk, 2010). By defining geopolitical enemies in strictly religious terms, these powers formed the basis of a civilizational division between the Christian and Islamic worlds. In both eras, religion did not guide politics, but rather, political ambitions were used to divide, categorize and weaponize believers against one another, permanently distorting the religious message of peace and substituting it with power.

This is evidenced by the fact that so many people have faith simply by virtue of birth; very few people are not born into a religion. As religious identity is so deeply woven into the fabric of human societies, it has always been the most potent mobilizing force available to leaders. The modern iteration of this instrumentalization is unique in the sense that it makes use of unprecedented tools of mass manipulation, leveraging mass media and globalized nationalist rhetoric to weaponize faith on a scale never seen or done before.

These dangers are thoroughly amplified when prophecy and the obsession with the messianic figure are used for political ends. Right-wing populists throughout the world use this rhetoric to portray themselves as saviors, and they more than most others, find that the only way to hold onto this newfound power is to rule by hate. Therefore, they pit groups against each other to maintain their own power, thus distorting religion to serve their own ends.

\vspace{2em}

In recent times, extremist Israeli leaders, have increasingly stated that their measures in Gaza and erstwhile Palestine, which the UN recognizes as an act of genocide, are hastening the arrival of their "Mahdi" or Messiah (Al Jazeera, 2024; Middle East Eye, 2024). This frames Israel's illegal wars in the Middle East in a religious light when in reality, religion is just the tool to justify these actions while the end is undeniably political: to gain more control and expand territory. Similarly, a US president who faced an assassination attempt recently attributed his survival directly to being chosen by God (Associated Press, 2024). He then used this claim of divine intervention as a justification to wage his own metaphorical and literal wars, asserting a mandate to "correct" the world in his own image because he had been rescued by God (Reich, 2026; Reuters, 2026). In both instances, religion is merely the instrument; the true end goal is to bring about political dominion. 

While exploring these themes, I was increasingly thinking about Frank Herbert's Dune: Messiah, a sci-fi fantasy which perils the danger of the messianic figure and holy wars. It details how destructive faith can be in the hands of someone who commands absolute power with it. The prophetic spice, melange, serves many allegories, one of which is blind faith: "It extends life and allows the adept to see his future, but it ties him to a cruel addiction and marks his eyes as yours are marked: total blue without any white. Your eyes, your organs of sight become one thing without any contrast, a single view" (Herbert, 1969). It is also eerily striking how the jihad detailed in the Dune chronology was religious in nature, waged at the behest of the messianic religious figure, but waged to retain control of this spice which also serves as the metaphor of oil in that reality, and by its end most of life was brought under a uniform government operated by the messianic figure. Therefore the instrument was religious, holy war. The end was political, to retake and keep control of the spice (oil) and the Universe.

It would be safe to say that the real world mirrors the world of Dune more than it does the introspective philosophy also associated with religion when it is not actively being used to wage division and conflict. One such reparative approach seen in the course is Ibn Sina's Treatise On Love. I want to explore how we can move from this bleak outlook on religion to one of healing, dignity and forbearance.

\vspace{2em}

In stark contrast to how religion and religious identity are used to foment division in modern geopolitics, Ibn Sina argues that the true nature of all beings is an innate, natural drive toward perfection, which he identifies as the "Pure Good" (Ibn Sina, 1945, p. 212). In this framework, religion is not a mechanism for conquest, but rather a disposition where the human soul seeks to assimilate to the Absolute Good through "noble deeds" and actions that are "in harmony with their nature" (Ibn Sina, 1945, p. 224). Rather than demanding blind allegiance, this philosophical approach to religion uplifts the soul by detailing how we can achieve divine love, a state of such fulfillment and annihilation that all earthly worries cease to be, and we lose ourselves in love of the Necessarily Existent, rather than in the lust for power. Ibn Sina suggests that to obtain the greatest love of the creator, one must imitate His qualities, namely, Justice and Intelligence. As such we can naturally infer that instrumentalizing God's word for worldly gain is neither just, nor intelligent. True spiritual elevation is brought about precisely through the exercise of "justice and intelligence" (Ibn Sina, 1945, p. 227).

Furthermore, Ibn Sina also draws a distinction between divine authority and the power wielded by earthly leaders. He states that the "Exalted King" (God) desires His creation to imitate His goodness and generously allows them to pursue this spiritual perfection (Ibn Sina, 1945, pp. 227-228). This stands in direct, condemning opposition to "earthly kings" who guard their power, becoming angry when others attempt to share in their authority, and actively preventing their subjects from achieving true nobility (Ibn Sina, 1945, p. 228). By reframing religion through Ibn Sina's lens, we strip away the worldly utility of holy war and political coercion. Faith is no longer an addiction of a single, blinding political view, and instead becomes an individual, rational pursuit of harmony, empathy, and universal love.

\vspace{2em}

To sum up, how can we move away from the dystopian use of religion to its place in the world as a paradigm of peace? To counter the instrumentalization of religion for war, it is vital to juxtapose it with a completely opposite outlook. The Treatise on Love frames Islam as a religion of profound love, centered on the contemplation of the divine and the rich traditions of Sufi philosophy, of self-annihilation and assimilation with the God to attain a path to peace in the world and the next.

In this framework, the purpose of faith is not to conquer territory or justify political domination, but to elevate the soul through an understanding of universal love, and beauty. It is deeply important to reconcile this aspect of religion such that it replaces its divisive misuse in the real world. The contrast highlights the tragedy of modern religious identity. The unifying, and healing aspects of the faith have been overshadowed by its utility as a weapon, for nation states and organizations to use as they see fit.

\vspace{2em}

\section*{References}

% This code creates a hanging indent layout perfect for APA format
\begin{list}{}{
  \setlength{\leftmargin}{0.5in}
  \setlength{\itemindent}{-0.5in}
  \setlength{\parsep}{\parskip}
}
\item Nauman Naqvi. (2024-). \textit{The collapse of the transcendant into the imminent}. Spoken in several courses throughout several semesters, (I first heard it in 2024).

\item Al Jazeera. (2024). \textit{Why the US and Israel are framing the ongoing conflict as a religious war?}. https://aje.news/8zl49x

\item Associated Press. (2024, July 14). \textit{Trump says `God alone prevented the unthinkable' after rally shooting}. 

\item Herbert, F. (1969). Excerpts from a death cell interview. In \textit{Dune Messiah}. G.P. Putnam's Sons.

\item Ibn Sina. (1945). A treatise on love (Ris\={a}lah f\={\i} al-'Ishq) (E. L. Fackenheim, Trans.). \textit{Mediaeval Studies}, (from the Coursepack).

\item Falk, A. (2010). Franks and Saracens: Reality and fantasy in the Crusades. Karnac Books.

\item Middle East Eye. (2024). \textit{Netanyahu: Israel will reach kingdom and make it Messiah's return}. https://www.middleeasteye.net/live-blog/live-blog-update/netanyahu-israel-will-reach-kingdom-and-make-it-messiahs-return

\item Reich, R. (2026). \textit{King Trump I versus Pope Leo XIV}. Robert Reich Substack. https://robertreich.substack.com/p/king-trump-i-versus-pope-leo-xiv

\item Reuters. (2026, April 5). \textit{Trump invokes religious rhetoric to praise Iran rescue drawing criticism}. https://www.reuters.com/business/media-telecom/trump-invokes-religious-rhetoric-praise-iran-rescue-drawing-criticism-2026-04-05/

\item Associated Press of Pakistan. (n.d.). \textit{Karbala tragedy: A timeless battle against tyranny and falsehood}. https://www.app.com.pk/domestic/karbala-tragedy-a-timeless-battle-against-tyranny-and-falsehood/

\end{list}

\end{document}
